\section{Исследование и построение решения задачи}
\label{sec:solution}

Построена математическая модель омни-платформы, формализован
принцип MPPI-управления и сформирована составная функция потерь,
отражающая прикладные требования из главы~\ref{sec:problem}.

\subsection{Кинематическая схема омни-платформы}

Робот представляет собой платформу с четырьмя колёсами типа Mecanum,
расположенными по углам прямоугольника со сторонами $2L$ (вдоль оси
$x_B$) и $2l$ (вдоль оси $y_B$) (рисунок~\ref{fig:omni_scheme}).
Ориентация роликов колёс составляет угол $\pm 45\degree{}$ к продольной оси платформы и
зеркально противоположна у соседних колёс.

\reportfig[0.7]{omni_scheme.png}{Кинематическая схема омни-платформы с колёсами Mecanum}{fig:omni_scheme}

Обозначения, используемые далее (в соответствии с
рисунком~\ref{fig:omni_scheme}):
\begin{itemize}
    \item $G$~--- геометрический центр платформы, совпадающий с центром
          связанной системы координат $x_B y_B$;
    \item $\psi$~--- угол ориентации корпуса относительно оси $x$
          глобальной системы координат;
    \item $L$~--- половина колёсной базы (расстояние от центра до оси
          колеса по оси $x_B$);
    \item $l$~--- половина колёсной колеи (расстояние от центра до оси
          колеса по оси $y_B$);
    \item $\dot{\varphi}_i$, $i = 1,\dots,4$~--- угловые скорости
          соответствующих колёс;
    \item $r$~--- радиус колеса.
\end{itemize}

Принцип работы колеса Mecanum показан на
рисунке~\ref{fig:omni_principle}: за счёт расположения свободно
вращающихся роликов под углом $45\degree{}$ к оси вращения колеса
сила реакции опоры создаёт составляющую как в продольном, так и в
поперечном направлении, что позволяет платформе с четырьмя такими
колёсами совершать голономное перемещение в плоскости~\cite{MorenoCaireta2020MPC}.

\reportfig[0.65]{omni wheel principle.png}{Принцип действия колеса Mecanum и формирование результирующей реакции опоры}{fig:omni_principle}

\subsection{Кинематика в непрерывном времени}

Пусть $v_x$, $v_y$~--- продольная и поперечная линейные скорости
корпуса в связанной системе координат, $\omega = \dot{\psi}$~--- угловая
скорость. Полный вектор обобщённых скоростей корпуса имеет вид
$\mathbf{u} = \begin{bmatrix} v_x & v_y & \omega \end{bmatrix}^{\top}$.

Связь между скоростью в связанной системе и производной глобальных
координат задаётся поворотом на угол $\psi$:
\begin{equation}
\dot{x} = v_x \cos\psi - v_y \sin\psi,
\label{eq:kin_x}
\end{equation}
\begin{equation}
\dot{y} = v_x \sin\psi + v_y \cos\psi,
\label{eq:kin_y}
\end{equation}
\begin{equation}
\dot{\psi} = \omega.
\label{eq:kin_psi}
\end{equation}

В матричной форме
\begin{equation}
\dot{\mathbf{x}} = R(\psi)\,\mathbf{u},\qquad
R(\psi) =
\begin{bmatrix}
\cos\psi & -\sin\psi & 0\\
\sin\psi & \cos\psi & 0\\
0 & 0 & 1
\end{bmatrix}.
\label{eq:kin_matrix}
\end{equation}

\subsection{Связь с угловыми скоростями колёс}

Для четырёхколёсной омни-платформы с радиусом колеса $r$, полубазой $L$
и полуколеёй $l$ прямая кинематика связывает обобщённые скорости
корпуса $\mathbf{u}$ с угловыми скоростями колёс $\dot{\varphi}_i$
следующим образом~\cite{MorenoCaireta2020MPC}:
\begin{equation}
\begin{bmatrix}
\dot{\varphi}_1\\[2pt]
\dot{\varphi}_2\\[2pt]
\dot{\varphi}_3\\[2pt]
\dot{\varphi}_4
\end{bmatrix}
=
\frac{1}{r}
\begin{bmatrix}
1 & -1 & -(L+l)\\
1 &  1 &  (L+l)\\
1 &  1 & -(L+l)\\
1 & -1 &  (L+l)
\end{bmatrix}
\begin{bmatrix}
v_x\\ v_y\\ \omega
\end{bmatrix},
\label{eq:jacobian}
\end{equation}
где строки матрицы соответствуют колёсам 1--4 в нумерации
рисунка~\ref{fig:omni_scheme} с учётом знаков ориентации роликов
($+45\degree{}$ для колёс 1 и 4, $-45\degree{}$ для колёс 2 и 3).

Обратное соотношение, используемое для проверки достижимости
управления, записывается как
\begin{equation}
\begin{bmatrix} v_x\\ v_y\\ \omega \end{bmatrix}
= \frac{r}{4}
\begin{bmatrix}
1 & 1 & 1 & 1\\
-1 & 1 & 1 & -1\\
-\frac{1}{L+l} & \frac{1}{L+l} & -\frac{1}{L+l} & \frac{1}{L+l}
\end{bmatrix}
\begin{bmatrix}
\dot{\varphi}_1\\ \dot{\varphi}_2\\ \dot{\varphi}_3\\ \dot{\varphi}_4
\end{bmatrix}.
\label{eq:inverse_jacobian}
\end{equation}

\subsection{Дискретная модель}

Для цифровой реализации контроллера непрерывная модель дискретизируется
с шагом $\Delta t$ методом Эйлера:
\begin{equation}
x_{k+1} = x_k + \Delta t\,\bigl(v_{x,k}\cos\psi_k - v_{y,k}\sin\psi_k\bigr),
\label{eq:disc_x}
\end{equation}
\begin{equation}
y_{k+1} = y_k + \Delta t\,\bigl(v_{x,k}\sin\psi_k + v_{y,k}\cos\psi_k\bigr),
\label{eq:disc_y}
\end{equation}
\begin{equation}
\psi_{k+1} = \psi_k + \Delta t\,\omega_k.
\label{eq:disc_psi}
\end{equation}

В компактной форме
\begin{equation}
\mathbf{x}_{k+1} = f(\mathbf{x}_k, \mathbf{u}_k).
\label{eq:disc_model}
\end{equation}

\subsection{Ограничения}

Физические ограничения на скорости и ускорения корпуса имеют вид
\begin{equation}
|v_{x,k}| \le v_x^{\max},\quad |v_{y,k}| \le v_y^{\max},\quad |\omega_k| \le \omega^{\max},
\label{eq:vel_limits}
\end{equation}
\begin{equation}
|\dot{v}_{x,k}| \le a_x^{\max},\quad |\dot{v}_{y,k}| \le a_y^{\max},\quad |\dot{\omega}_k| \le \alpha^{\max}.
\label{eq:acc_limits}
\end{equation}

Условие безопасности относительно препятствий записывается как
\begin{equation}
d_{i,k} = \|\mathbf{p}_k - \mathbf{p}^{\text{obs}}_{i,k}\|_2 \ge d_{\min},
\quad i = 1,\dots,n_{\text{obs}}.
\label{eq:safety_constraint}
\end{equation}

\subsection{Принцип модельного прогнозирующего управления}

На шаге $k$ модельное прогнозирующее управление (MPC) формирует
последовательность управляющих воздействий
\begin{equation}
\mathbf{U}_k^{\star} =
\left\{
\mathbf{u}_{k|k}^{\star},
\mathbf{u}_{k+1|k}^{\star},
\dots,
\mathbf{u}_{k+N-1|k}^{\star}
\right\},
\label{eq:mpc_sequence}
\end{equation}
где $N$~--- горизонт прогнозирования. Прогноз состояния строится
рекурсивно с использованием модели~\eqref{eq:disc_model}:
\begin{equation}
\mathbf{x}_{k+j+1|k} = f(\mathbf{x}_{k+j|k}, \mathbf{u}_{k+j|k}),
\quad j = 0,\dots,N-1.
\label{eq:mpc_rollout}
\end{equation}

Оптимизируемый функционал представляет сумму этапных и терминальных
затрат:
\begin{equation}
J_k = \sum_{j=0}^{N-1}
\ell\bigl(\mathbf{x}_{k+j|k}, \mathbf{u}_{k+j|k}\bigr)
+ \ell_f\bigl(\mathbf{x}_{k+N|k}\bigr).
\label{eq:mpc_cost}
\end{equation}

По принципу receding horizon применяется только первое управление
$\mathbf{u}_k = \mathbf{u}_{k|k}^{\star}$, после чего горизонт сдвигается,
измеряется новое состояние и задача оптимизации решается заново.

\subsection{Стохастическая реализация MPPI}

В работе используется стохастическая реализация MPC --- MPPI, имеющаяся
в составе пакета \texttt{nav2\_mppi\_controller} навигационного стека
Nav2~\cite{Alkhaddad2023MPC}. Для каждого сэмпла $m = 1,\dots,M$
генерируется возмущение управления
\begin{equation}
\boldsymbol{\epsilon}_{j}^{(m)} \sim \mathcal{N}(0,\Sigma),
\quad m = 1,\dots,M,\quad j = 0,\dots,N-1,
\label{eq:mppi_noise}
\end{equation}
и формируется траектория $m$-го сэмпла с управлением
\begin{equation}
\mathbf{u}_{k+j}^{(m)} = \bar{\mathbf{u}}_{k+j} + \boldsymbol{\epsilon}_{j}^{(m)},
\label{eq:mppi_u}
\end{equation}
где $\bar{\mathbf{u}}_{k+j}$~--- номинальное управление на шаге $j$.

Суммарная стоимость $m$-й траектории
\begin{equation}
S^{(m)} = \sum_{j=0}^{N-1}
\ell\bigl(\mathbf{x}_{k+j}^{(m)}, \mathbf{u}_{k+j}^{(m)}\bigr)
+ \ell_f\bigl(\mathbf{x}_{k+N}^{(m)}\bigr).
\label{eq:mppi_S}
\end{equation}

Вес траектории формируется экспоненциальным правилом
\begin{equation}
w^{(m)} =
\frac{\exp\bigl(-\tfrac{1}{\lambda} S^{(m)}\bigr)}
     {\sum\limits_{q=1}^{M} \exp\bigl(-\tfrac{1}{\lambda} S^{(q)}\bigr)},
\label{eq:mppi_weights}
\end{equation}
где $\lambda > 0$~--- температурный параметр (в реализации Nav2~MPPI
обозначается \texttt{temperature}).

Обновление номинального управления происходит по взвешенному среднему
возмущений:
\begin{equation}
\bar{\mathbf{u}}_{k+j}^{\text{new}} =
\bar{\mathbf{u}}_{k+j} + \sum_{m=1}^{M} w^{(m)}\,\boldsymbol{\epsilon}_{j}^{(m)}.
\label{eq:mppi_update}
\end{equation}

\subsection{Архитектура функции потерь}

Согласно требованиям задачи (глава~\ref{sec:problem}) функция потерь
должна одновременно обеспечивать:
\begin{itemize}
    \item поддержание целевой дистанции до пользователя;
    \item центрирование цели в поле зрения камеры по горизонтали;
    \item безопасность относительно препятствий;
    \item плавность управляющих воздействий;
    \item корректное поведение при потере наблюдаемости цели.
\end{itemize}

Введём следующие отклонения на шаге $k$:
\begin{equation}
e^{d}_{k} = \|\mathbf{p}_k - \mathbf{x}^{t}_k\|_2 - d^{\star},
\label{eq:err_dist}
\end{equation}
\begin{equation}
e^{\phi}_{k} =
\operatorname{atan2}\bigl(y^{t}_k - y_k,\; x^{t}_k - x_k\bigr) - \psi_k,
\label{eq:err_bearing}
\end{equation}
где $d^{\star}$~--- желаемая дистанция, а $e^{\phi}_{k}$ приводится к
$(-\pi,\pi]$.

Этапная функция потерь принимается в виде взвешенной суммы
\begin{align}
      \ell_k = &\; w_d\,(e^{d}_{k})^{2} + w_{\phi}\,(e^{\phi}_{k})^{2} + w_{u}\,\|\mathbf{u}_k\|_2^{2} \nonumber \\
      &+ w_{\Delta u}\,\|\mathbf{u}_k - \mathbf{u}_{k-1}\|_2^{2} + w_{\text{obs}}\,\Psi_{\text{obs},k} \nonumber \\
      &+ w_{\text{pf}}\,\Psi_{\text{pf},k} + w_{\text{hold}}\,\Psi_{\text{hold},k},
      \label{eq:stage_cost}
      \end{align}
где
\begin{equation}
\Psi_{\text{obs},k} =
\sum_{i=1}^{n_{\text{obs}}}
\frac{1}{(d_{i,k} - d_{\text{safe}})^{2} + \varepsilon},
\quad \varepsilon > 0,
\label{eq:psi_obs}
\end{equation}
\begin{equation}
\Psi_{\text{hold},k} = (1 - \nu_k)\,\|\mathbf{u}_k\|_2^{2},
\label{eq:psi_hold}
\end{equation}
а $\nu_k \in \{0,1\}$~--- бинарный индикатор видимости цели.
Слагаемое $\Psi_{\text{pf},k}$ задаёт направленный (анизотропный) штраф
вдоль градиента поля препятствий и определено в
подразделе~\ref{sec:pf_cost}.
Терминальная стоимость имеет вид
\begin{equation}
\ell_f =
w_f^{d}\,(e^{d}_{k+N})^{2}
+ w_f^{\phi}\,(e^{\phi}_{k+N})^{2}.
\label{eq:terminal_cost}
\end{equation}

\subsection{Направленный штраф поля препятствий}
\label{sec:pf_cost}

Изотропный барьер~\eqref{eq:psi_obs} штрафует только сам факт
близости препятствия и не различает направления движения, что для
голономной омни-платформы (свободный выбор $v_x$, $v_y$) недостаточно:
сэмплы, идущие \emph{вдоль} края препятствия, и сэмплы, направленные
\emph{в} препятствие, получают одинаковый штраф. Для устранения этого
эффекта в функцию потерь дополнительно вводится направленный (потенциального
поля) штраф $\Psi_{\text{pf},k}$.

Пусть $c(x,y) \in [0,1]$~--- нормированное на максимум значение
costmap в точке $(x,y)$, $\nabla c$~--- его конечно-разностный градиент.
Введём внутреннее расстояние до препятствия $\tilde{d}(c)$,
восстанавливаемое из инфляционной модели costmap. Тогда
\begin{equation}
\hat{\mathbf{f}}_k =
-\frac{\nabla c(\mathbf{p}_k)}{\|\nabla c(\mathbf{p}_k)\|_2}
\label{eq:pf_dir}
\end{equation}
~--- единичное направление отталкивания от препятствия в точке
$\mathbf{p}_k$.

Слагаемое $\Psi_{\text{pf},k}$ строится как сумма двух
неотрицательных составляющих: проксимальной и направленной.
Проксимальная составляющая активна, когда восстановленное расстояние
до препятствия меньше $d^{\star}_{\text{pf}}$:
\begin{equation}
\rho_k = \max\bigl(0,\; d^{\star}_{\text{pf}} - \tilde{d}(c(\mathbf{p}_k))\bigr).
\label{eq:pf_prox}
\end{equation}
Направленная составляющая штрафует движение
\emph{в~сторону} препятствия~--- проекцию единичного вектора
скорости $\hat{\mathbf{v}}_k = \mathbf{v}_k / \|\mathbf{v}_k\|_2$ на
направление, противоположное отталкиванию:
\begin{equation}
\sigma_k = \max\bigl(0,\; -\langle \hat{\mathbf{v}}_k,\, \hat{\mathbf{f}}_k \rangle\bigr)
\,(1 + \rho_k).
\label{eq:pf_dir_score}
\end{equation}
Полное слагаемое имеет вид
\begin{equation}
\Psi_{\text{pf},k} = \mathbf{1}\!\bigl[c(\mathbf{p}_k) \ge c_{\text{act}}\bigr]\,
\bigl(w_{\rho}\,\rho_k + w_{\sigma}\,\sigma_k\bigr),
\label{eq:psi_pf}
\end{equation}
где $\mathbf{1}[\cdot]$~--- индикатор активации по порогу costmap
$c_{\text{act}}$, а $w_{\rho}$, $w_{\sigma}$~--- внутренние веса
проксимальной и направленной составляющих.

Активация по порогу $c_{\text{act}}$ выключает штраф вдали от
препятствий, что делает его локальным регуляризатором поведения и
не искажает преследование цели в свободном пространстве. На
расстоянии до текущей цели меньше $d_{\text{ng}}$ слагаемое также
обнуляется, чтобы исключить конфликт с подходом к финальной
позе~\cite{Alkhaddad2023MPC}.

\subsection{Интерпретация весовых коэффициентов}

Веса $w_{(\cdot)}$ регулируют относительную значимость слагаемых
функции потерь. Характерные соотношения:
\begin{itemize}
    \item $w_{\text{obs}}, w_{\text{pf}} \gg w_{d}, w_{\phi}$~---
          приоритет безопасности над точностью сопровождения;
    \item увеличение $w_{\Delta u}$ уменьшает рывки управления:
          $\|\mathbf{u}_k - \mathbf{u}_{k-1}\|_2 \to 0$;
    \item $w_{\text{hold}} \gg 0$ при $\nu_k = 0$ заставляет робот
          остановиться: $\mathbf{u}_k \to \mathbf{0}$;
    \item внутренний $w_{\sigma}$ в~\eqref{eq:pf_dir_score} подавляет
          движение в препятствие на голономных направлениях, не
          мешая движению вдоль границы препятствия.
\end{itemize}

\subsection{Декомпозиция на критики Nav2/MPPI}

Любое слагаемое функции потерь в \eqref{eq:stage_cost} представляется
независимым модулем~--- критиком траектории. Формально
\begin{equation}
\ell_k = \sum_{r=1}^{R} \ell_k^{(r)},
\label{eq:decomp_stage}
\end{equation}
\begin{equation}
S^{(m)} = \sum_{j=0}^{N-1} \sum_{r=1}^{R} \ell_{k+j}^{(r,m)},
\label{eq:decomp_S}
\end{equation}
где $\ell_k^{(r)}$~--- вклад $r$-го критика. Такая декомпозиция
обеспечивает:
\begin{itemize}
    \item независимую параметрическую настройку каждого требования;
    \item интерпретируемость поведения в терминах отдельных целей;
    \item модульное расширение системы при добавлении новых требований
          к управлению.
\end{itemize}

Соответствие слагаемых функции~\eqref{eq:stage_cost} критикам
реализации Nav2~MPPI приведено в главе~\ref{sec:implementation}.

\subsection{Контур восприятия цели}

Для связи функции потерь с измерениями сенсоров необходимо пересчитать
положение обнаруженной цели в то же координатное пространство, в котором
формируются прогнозы траектории. Схема пересчёта включает следующие
этапы:
\begin{enumerate}
    \item детектирование цели по цветовой сегментации изображения с
          фронтальной камеры;
    \item вычисление углового направления на цель по положению центра
          ограничивающего прямоугольника и известному горизонтальному
          углу поля зрения камеры;
    \item оценка дальности до цели по данным 2D-лидара в окрестности
          полученного направления (медиана допустимых значений в
          заданном окне индексов);
    \item преобразование точки из связанной с роботом системы координат
          в глобальную карту через дерево преобразований TF.
\end{enumerate}

Такая схема обеспечивает согласованность математической модели и
измерений, а также позволяет использовать один и тот же кадр карты
для всех слагаемых функции потерь.
